% ============================================================================
% Gneiting, T. & Raftery, A. E. (2007). Strictly Proper Scoring Rules,
% Prediction, and Estimation. JASA 102(477), 359--378. doi:10.1198/016214506000001437
%
% PARTIAL TRANSCRIPTION (no official LaTeX source exists). Transcribed by
% Claude from paper.pdf; equation numbers match the published paper.
% Included: Sec 1 (optimum score estimation idea), 2.1, 4.2, 4.3, 5.1, 5.2.
% Omitted: categorical scores (Sec 3), density scores (4.1), 4.4, 5.3,
%          quantile/interval scores (6), Bayes factors/CV (7), case study (8),
%          optimum score estimation details (9), future work (10).
% Verify any equation you rely on against paper.pdf.
%
% SIGN CONVENTION: this paper uses POSITIVELY oriented scores (higher = better,
% forecaster maximises). CRPS as defined in (20)/(21) is <= 0. The usual ML loss
% is the negatively oriented CRPS^* = -CRPS (see end of Sec 4.2).
% ============================================================================
\documentclass{article}
\usepackage{amsmath,amssymb,amsthm}
\newtheorem{theorem}{Theorem}
\newtheorem*{assumption}{Assumption 1}
\newcommand{\1}{\mathbb{1}}
\newcommand{\R}{\mathbb{R}}
\newcommand{\bx}{\mathbf{x}}
\newcommand{\bX}{\mathbf{X}}
\newcommand{\by}{\mathbf{y}}
\newcommand{\bY}{\mathbf{Y}}
\begin{document}

\section*{1. Introduction (excerpt): optimum score estimation}
Suppose we wish to fit a parametric model $P_\theta$ based on a sample
$X_1,\dots,X_n$. We might measure goodness-of-fit by the mean score
\[
  \mathcal{S}_n(\theta) = \frac{1}{n}\sum_{i=1}^n S(P_\theta, X_i),
\]
where $S$ is a strictly proper scoring rule. If $\theta_0$ is the true parameter
value, asymptotic arguments indicate that
$\arg\max_\theta \mathcal{S}_n(\theta) \to \theta_0$ as $n\to\infty$. The
\emph{optimum score estimator} is $\hat\theta_n = \arg\max_\theta \mathcal{S}_n(\theta)$
(minimum contrast estimation; a special case of $M$-estimation; maximum
likelihood is the special case of the logarithmic score).

\section*{2.1 Proper scoring rules and convex functions}
Let $\Omega$ be a sample space, $\mathcal{A}$ a $\sigma$-algebra of subsets of
$\Omega$, and $\mathcal{P}$ a convex class of probability measures on
$(\Omega,\mathcal{A})$. A \emph{probabilistic forecast} is any $P\in\mathcal{P}$.
A \emph{scoring rule} is any extended real-valued function
$S:\mathcal{P}\times\Omega\to\overline{\R}$ such that $S(P,\cdot)$ is
$\mathcal{P}$-quasi-integrable for all $P\in\mathcal{P}$. If the forecast is $P$
and $\omega$ materialises, the forecaster's reward is $S(P,\omega)$. Write
\[
  S(P,Q) = \int S(P,\omega)\,\mathrm{d}Q(\omega)
\]
for the expected score under $Q$ when the forecast is $P$. $S$ is \emph{proper}
relative to $\mathcal{P}$ if
\begin{equation}\tag{1}
  S(Q,Q) \ge S(P,Q) \quad \text{for all } P,Q\in\mathcal{P}.
\end{equation}
It is \emph{strictly proper} if (1) holds with equality iff $P=Q$. If $S$ is
proper, $c>0$ a constant and $h$ a $\mathcal{P}$-integrable function, then
\begin{equation}\tag{2}
  S^*(P,\omega) = c\,S(P,\omega) + h(\omega)
\end{equation}
is also proper (strictly proper if $S$ is); $S$ and $S^*$ are \emph{equivalent},
and \emph{strongly equivalent} if $c=1$.

\section*{4.2 Continuous ranked probability score}
Let $\mathcal{P}$ be the Borel probability measures on $\R$; identify a forecast
with its CDF $F$. The CRPS is
\begin{equation}\tag{20}
  \mathrm{CRPS}(F,x) = -\int_{-\infty}^{\infty} \bigl(F(y) - \1\{y\ge x\}\bigr)^2\,\mathrm{d}y ,
\end{equation}
the integral of the Brier scores for the binary forecasts at all real thresholds.
By lemma 2.2 of Baringhaus \& Franz (2004) or identity (17) of Sz\'ekely \&
Rizzo (2005),
\begin{equation}\tag{21}
  \mathrm{CRPS}(F,x) = \tfrac12 E_F|X-X'| - E_F|X-x| ,
\end{equation}
where $X,X'$ are independent copies of a random variable with CDF $F$ and finite
first moment. For a Gaussian predictive distribution,
\[
  \mathrm{CRPS}\bigl(\mathcal{N}(\mu,\sigma^2),x\bigr)
  = \sigma\left[\frac{1}{\sqrt\pi} - 2\varphi\!\left(\frac{x-\mu}{\sigma}\right)
  - \frac{x-\mu}{\sigma}\left(2\Phi\!\left(\frac{x-\mu}{\sigma}\right)-1\right)\right],
\]
with $\varphi,\Phi$ the standard normal PDF and CDF. If the predictive
distribution is a sample of size $n$, the right side of (20) can be evaluated
via order statistics in $\mathcal{O}(n\log n)$ operations (Hersbach 2000, sec.~4.b).

The CRPS is proper relative to $\mathcal{P}$ and \textbf{strictly proper
relative to the subclass $\mathcal{P}_1$ of Borel measures with finite first
moment}. The associated expected score (generalised entropy) function is
\[
  G(F) = -\int_{-\infty}^{\infty} F(y)\bigl(1-F(y)\bigr)\,\mathrm{d}y = -\tfrac12 E_F|X-X'| ,
\]
and the divergence function is
\[
  d(F,G) = \int_{-\infty}^{\infty} \bigl(F(y)-G(y)\bigr)^2\,\mathrm{d}y ,
\]
symmetric and of Cram\'er--von Mises type.

It is typically used in negative orientation, $\mathrm{CRPS}^*(F,x) = -\mathrm{CRPS}(F,x)$,
so (21) becomes
\[
  \mathrm{CRPS}^*(F,x) = E_F|X-x| - \tfrac12 E_F|X-X'| .
\]
In negative orientation it has the units of the observations and reduces to the
absolute error when $F$ is a point measure (deterministic forecast).

\section*{4.3 Energy score}
Let $\mathcal{P}_\beta$, $\beta\in(0,2)$, be the Borel probability measures $P$
on $\R^m$ with $E_P\|\bX\|^\beta<\infty$ ($\|\cdot\|$ Euclidean). The
\emph{energy score} is
\begin{equation}\tag{22}
  \mathrm{ES}(P,\bx) = \tfrac12 E_P\|\bX-\bX'\|^\beta - E_P\|\bX-\bx\|^\beta ,
\end{equation}
with $\bX,\bX'$ independent copies with distribution $P\in\mathcal{P}_\beta$.
It reduces to the CRPS (21) when $\beta=1$, $m=1$. Theorem 1 of Sz\'ekely (2003):
the energy score is strictly proper relative to $\mathcal{P}_\beta$. In the
limiting case $\beta=2$,
\begin{equation}\tag{23}
  \mathrm{ES}(P,\bx) = -\|\boldsymbol\mu_P - \bx\|^2 ,
\end{equation}
where $\boldsymbol\mu_P$ is the mean vector of $P$: regular and proper but
\textbf{not strictly proper} relative to $\mathcal{P}_2$.

For all Borel probability measures on $\R^m$ and $\beta\in(0,2)$,
\begin{equation}\tag{24}
  \mathrm{ES}(P,\bx) = -\frac{\beta\,2^{\beta-2}\,\Gamma\!\left(\frac m2+\frac\beta2\right)}
  {\pi^{m/2}\,\Gamma\!\left(1-\frac\beta2\right)}
  \int_{\R^m} \frac{\bigl|\phi_P(\by) - e^{i\langle \bx,\by\rangle}\bigr|^2}{\|\by\|^{m+\beta}}\,\mathrm{d}\by ,
\end{equation}
where $\phi_P$ is the characteristic function of $P$. If $P\in\mathcal{P}_\beta$
the right sides of (22) and (24) agree: the score is a weighted distance between
the characteristic function of $P$ and that of the point measure at the outcome.

\section*{5.1 Kernel scores}
A real-valued $g$ on $\Omega\times\Omega$ is a \emph{negative definite kernel}
if it is symmetric and $\sum_{i=1}^n\sum_{j=1}^n a_i a_j g(x_i,x_j)\le 0$ for all
positive integers $n$, all $a_1,\dots,a_n\in\R$ summing to $0$, and all
$x_1,\dots,x_n\in\Omega$.

\begin{theorem}[Theorem 4]
Let $\Omega$ be a Hausdorff space and $g$ a nonnegative, continuous negative
definite kernel on $\Omega\times\Omega$. For a Borel probability measure $P$ on
$\Omega$, let $X,X'$ be independent with distribution $P$. Then
\begin{equation}\tag{28}
  S(P,x) = \tfrac12 E_P\, g(X,X') - E_P\, g(X,x)
\end{equation}
is proper relative to the class of Borel probability measures $P$ on $\Omega$
for which $E_P\, g(X,X')$ is finite.
\end{theorem}
\begin{proof}
Let $X,X'$ and $Y,Y'$ be independent with distributions $P$ and $Q$. We need
\begin{equation}\tag{29}
  -\tfrac12 E_Q\, g(Y,Y') \ge \tfrac12 E_P\, g(X,X') - E_{P,Q}\, g(X,Y).
\end{equation}
If $E_{P,Q}\,g(X,Y)$ is infinite the inequality is trivial; if finite, theorem
2.1 of Berg et al.\ (1984, p.~235) implies (29).
\end{proof}

\noindent\textbf{Example 7 (Brier).} $\Omega=\{1,0\}$, $g(0,0)=g(1,1)=0$,
$g(0,1)=g(1,0)=1$: (28) recovers the quadratic/Brier score.

\noindent\textbf{Example 8 (CRPS).} $\Omega=\R$, $g(x,x')=|x-x'|$: (28) gives
the CRPS (21).

\noindent\textbf{Example 9 (Energy score).} $\Omega=\R^m$, $\beta\in(0,2)$,
$g(\bx,\bx')=\|\bx-\bx'\|^\beta$: (28) gives the energy score (22).

\noindent\textbf{Example 10 (circular CRPS).} $\Omega=\mathbb{S}$, angular
distance $\alpha(\theta,\theta')$ is negative definite (Gneiting 1998), so
\begin{equation}\tag{30}
  S(P,\theta) = \tfrac12 E_P\,\alpha(\Theta,\Theta') - E_P\,\alpha(\Theta,\theta)
\end{equation}
is proper.

For $\bx\in\R^m$ and $\alpha\in(0,\infty]$ define
$\|\bx\|_\alpha = (\sum_{i=1}^m|x_i|^\alpha)^{1/\alpha}$ ($\max_i|x_i|$ if
$\alpha=\infty$). Then for $\beta>0$ the kernel $g(\bx,\bx')=\|\bx-\bx'\|_\alpha^\beta$
is negative definite iff:

\begin{assumption}
(a) $m=1$, $\alpha\in(0,\infty]$, $\beta\in(0,2]$; or (b) $m\ge2$,
$\alpha\in(0,2]$, $\beta\in(0,\alpha]$; or (c) $m=2$, $\alpha\in(2,\infty]$,
$\beta\in(0,1]$.
\end{assumption}

\noindent\textbf{Example 11 (non-Euclidean energy score).} Under Assumption 1,
\[
  S(P,\bx) = \tfrac12 E_P\|\bX-\bX'\|_\alpha^\beta - E_P\|\bX-\bx\|_\alpha^\beta
\]
is proper relative to Borel measures on $\R^m$ with $E_P\|\bX-\bX'\|_\alpha^\beta<\infty$.
$m=1$ or $\alpha=2$ recovers the energy score. (Mattner 1997: if $\alpha\ge1$,
$E_{P,Q}\|\bX-\bY\|_\alpha^\beta$ finite iff $E_P\|\bX\|_\alpha^\beta$ and
$E_Q\|\bY\|_\alpha^\beta$ finite.)

\begin{theorem}[Theorem 5]
Let $\psi$ be continuous on $[0,\infty)$ with $-\psi'$ completely monotone and
not constant. For a Borel probability measure $P$ on $\R^m$ with $\bX,\bX'$
independent with distribution $P$,
\[
  S(P,\bx) = \tfrac12 E_P\,\psi(\|\bX-\bX'\|_2^2) - E_P\,\psi(\|\bX-\bx\|_2^2)
\]
is \textbf{strictly} proper relative to Borel $P$ on $\R^m$ with
$E_P\,\psi(\|\bX-\bX'\|_2^2)<\infty$. (Immediate from theorem 2.2 of Mattner
1997.) With $\psi(t)=t^{\beta/2}$, $\beta\in(0,2)$, this gives strict propriety
of the energy score for $P$ with $E_P\|\bX\|_2^\beta<\infty$.
\end{theorem}
($\eta$ on $(0,\infty)$ is completely monotone if $(-1)^k\eta^{(k)}(t)\ge0$ for
all $k\ge0$, $t>0$.)

\section*{5.2 Inequalities of Hoeffding type and positive definite kernels}
If $E_P\,g(X,X')$ and $E_Q\,g(Y,Y')$ are finite, (29) is the Hoeffding-type
inequality
\begin{equation}\tag{31}
  2E_{P,Q}\,g(X,Y) - E_P\,g(X,X') - E_Q\,g(Y,Y') \ge 0 .
\end{equation}
(Sz\'ekely \& Rizzo 2005, thm 1: nearly identical result and a converse; if $g$
is not negative definite there are counterexamples to (31) and the score is
improper.) If $\Omega$ is a group and $g(x,x')=g(-x,-x')$,
\begin{equation}\tag{32}
  E_P\,g(X,-X') \ge E_P\,g(X,X').
\end{equation}
For $\Omega=\R^m$ under Assumption 1, (31) and (32) become
\begin{equation}\tag{33}
  2E\|\bX-\bY\|_\alpha^\beta - E\|\bX-\bX'\|_\alpha^\beta - E\|\bY-\bY'\|_\alpha^\beta \ge 0
\end{equation}
and
\begin{equation}\tag{34}
  E\|\bX-\bX'\|_\alpha^\beta \le E\|\bX+\bX'\|_\alpha^\beta .
\end{equation}
[Note: the left side of (33) is the (generalised) \emph{energy distance}
between $P$ and $Q$.] For a group $\Omega$ with $g(x,x')=g(-x,-x')$,
\begin{equation}\tag{35}
  h(x,x') = g(x,-x') - g(x,x')
\end{equation}
is a positive definite kernel; in particular, under Assumption 1,
\begin{equation}\tag{36}
  h(\bx,\bx') = \|\bx+\bx'\|_\alpha^\beta - \|\bx-\bx'\|_\alpha^\beta
\end{equation}
is positive definite on $\R^m$.

\end{document}
